\documentclass[aspectratio=169,11pt]{beamer}

\usepackage[utf8]{inputenc}
\usepackage[T1]{fontenc}
\usepackage[brazil]{babel}
\usepackage{amsmath}
\usepackage{amssymb}
\usepackage{graphicx}
\usepackage{svg}
\usepackage{tikz}
\usepackage[most]{tcolorbox}
\usepackage{lmodern}
\usepackage{array}
\usepackage{tabularx}
\usepackage{booktabs}

\usetikzlibrary{shadows,shadows.blur,positioning,arrows.meta}
\svgsetup{inkscapelatex=false}

\definecolor{mainblue}{HTML}{3B35B9}
\definecolor{darkblue}{HTML}{16175A}
\definecolor{lightbox}{HTML}{ECECF7}
\definecolor{alertred}{HTML}{D60000}
\definecolor{alertbg}{HTML}{F8E5E5}

\setbeamertemplate{navigation symbols}{}
\setbeamercolor{background canvas}{bg=white}
\setbeamercolor{normal text}{fg=black}
\setbeamercolor{item}{fg=mainblue}
\setbeamerfont{frametitle}{size=\Large,series=\mdseries}
\setbeamerfont{normal text}{size=\normalsize}
\setbeamersize{text margin left=0.60cm,text margin right=0.60cm}
\setbeamertemplate{itemize item}{\color{mainblue}$\blacktriangleright$}
\setbeamertemplate{itemize subitem}{\color{mainblue}$\blacktriangleright$}
\setlength{\leftmargini}{1.5em}
\setlength{\leftmarginii}{1.3em}

\setbeamercolor{frametitle}{bg=mainblue,fg=white}
\setbeamertemplate{frametitle}{%
    \nointerlineskip%
    \begin{beamercolorbox}[wd=\paperwidth,ht=0.72cm,dp=0.17cm,leftskip=0.32cm]{frametitle}%
        \usebeamerfont{frametitle}\insertframetitle
    \end{beamercolorbox}%
}

\setbeamercolor{footerleft}{bg=darkblue,fg=white}
\setbeamercolor{footerright}{bg=mainblue,fg=white}
\newcommand{\footersafespace}{0.30cm}
\newcommand{\footerheight}{0.25cm}
\newcommand{\footerdepth}{0.09cm}
\setbeamertemplate{footline}{%
    \leavevmode%
    \vbox{%
        \vskip\footersafespace%
        \hbox to \paperwidth{%
            \begin{beamercolorbox}[wd=0.85\paperwidth,ht=\footerheight,dp=\footerdepth,leftskip=0.10cm]{footerleft}%
                \tiny INF01146 -- Avaliação de Desempenho -- Trabalho Final: Etapa 2%
            \end{beamercolorbox}%
            \begin{beamercolorbox}[wd=0.15\paperwidth,ht=\footerheight,dp=\footerdepth,rightskip=0.10cm]{footerright}%
                \hfill\tiny\insertframenumber/\inserttotalframenumber%
            \end{beamercolorbox}%
        }%
    }%
}

\tcbset{slidebox/.style={enhanced,boxrule=0pt,arc=2mm,outer arc=2mm,left=2.2mm,right=2.2mm,top=1.4mm,bottom=1.4mm,toptitle=1mm,bottomtitle=1mm,drop shadow={shadow xshift=1.4mm,shadow yshift=-1.4mm,opacity=0.35}}}
\newtcolorbox{questionbox}[1]{slidebox,colback=lightbox,colframe=lightbox,coltitle=white,colbacktitle=darkblue,title=#1,fonttitle=\normalsize}
\newtcolorbox{impactbox}[1]{slidebox,colback=alertbg,colframe=alertbg,coltitle=white,colbacktitle=alertred,title=#1,fonttitle=\normalsize}
\newcommand{\chartimage}[2][]{%
    \begin{center}\includesvg[width=\linewidth,height=0.68\textheight,keepaspectratio,#1]{#2}\end{center}%
}
\newcommand{\chartcaption}[1]{%
    \vspace{-0.16cm}\par\footnotesize\textcolor{darkblue}{\textbf{Leitura.}} #1\par%
}
\tikzset{
  flowbox/.style={rounded corners=2mm, fill=lightbox, draw=mainblue,
    line width=0.6pt, minimum height=1.25cm, align=center, text width=2.18cm,
    font=\small},
  flowarrow/.style={-{Stealth[length=2.0mm]}, mainblue, line width=0.7pt}
}

\title{Avaliação de Desempenho Computacional}
\subtitle{Algoritmo Genético -- Etapa 2}
\author{Felipe Avencourt Soares -- 342904\\[0.20cm]Leonardo Rocha dos Santos -- 341831}
\institute{Instituto de Informática -- UFRGS 2026/2}
\date{}

\begin{document}

{%
\setbeamertemplate{frametitle}{}%
\begin{frame}
    \vfill
    \begin{center}
        \begin{tcolorbox}[enhanced,width=0.92\textwidth,colback=mainblue,colframe=mainblue,arc=2mm,boxrule=0pt,left=3mm,right=3mm,top=5mm,bottom=5mm,drop shadow={shadow xshift=1.7mm,shadow yshift=-1.7mm,opacity=0.35}]
            \centering
            {\color{white}\LARGE\inserttitle\par}
            \vspace{0.15cm}
            {\color{white}\Large\insertsubtitle}
        \end{tcolorbox}
        \vspace{0.75cm}
        {\large\insertauthor\par}
        \vspace{0.35cm}
        {\small\insertinstitute\par}
    \end{center}
    \vfill
\end{frame}
}

\begin{frame}{Nossa jornada nesta apresentação}
    \begin{enumerate}
        \item A pergunta e o algoritmo que medimos.
        \item O experimento: fatores, controles, réplicas e métricas.
        \item A cadeia que transforma uma execução em evidência auditável.
        \item Os resultados: tempo, vazão, fases, contadores e variabilidade.
        \item O que já validamos, os limites e o próximo passo.
    \end{enumerate}
    \vspace{0.32cm}
    \begin{questionbox}{Ideia central}
        Nós estudamos o custo de executar o algoritmo. A qualidade das soluções
        encontradas pelo algoritmo fica fora deste experimento.
    \end{questionbox}
\end{frame}

\begin{frame}{Pergunta de pesquisa e recorte}
    \begin{questionbox}{Pergunta herdada da Etapa 1}
        Como o tamanho da população $N$ e o comprimento do cromossomo $L$
        afetam o desempenho computacional de um algoritmo genético sequencial?
    \end{questionbox}
    \vspace{0.28cm}
    \begin{columns}[T,onlytextwidth]
        \begin{column}{0.48\textwidth}
            \textbf{O que varia}
            \begin{itemize}
                \item $N$: número de indivíduos na população;
                \item $L$: número de genes em cada cromossomo.
            \end{itemize}
        \end{column}
        \begin{column}{0.48\textwidth}
            \textbf{O que permanece fixo}
            \begin{itemize}
                \item uma thread e $G=100$ gerações;
                \item regras do algoritmo e ambiente registrado.
            \end{itemize}
        \end{column}
    \end{columns}
    \vspace{0.25cm}
    \begin{impactbox}{Escopo}
        A resposta procurada é sobre tempo de execução e trabalho por gene.
        Este estudo não compara a qualidade de otimização do AG.
    \end{impactbox}
\end{frame}

\begin{frame}{O algoritmo genético em uma geração}
    \begin{center}
    \begin{tikzpicture}[node distance=0.36cm and 0.33cm]
        \node[flowbox] (population) {\textbf{População}\\$N$ cromossomos\\com $L$ genes};
        \node[flowbox, right=of population] (fitness) {\textbf{Avaliação}\\OneMax conta\\os genes iguais a 1};
        \node[flowbox, right=of fitness] (selection) {\textbf{Seleção}\\torneio escolhe\\os pais};
        \node[flowbox, right=of selection] (crossover) {\textbf{Crossover}\\um ponto combina\\os pais};
        \node[flowbox, below=of crossover] (mutation) {\textbf{Mutação}\\cada gene pode\\mudar de 0 para 1\\ou de 1 para 0};
        \node[flowbox, left=of mutation] (replacement) {\textbf{Substituição}\\a nova população\\entra no lugar da anterior};
        \draw[flowarrow] (population) -- (fitness);
        \draw[flowarrow] (fitness) -- (selection);
        \draw[flowarrow] (selection) -- (crossover);
        \draw[flowarrow] (crossover) -- (mutation);
        \draw[flowarrow] (mutation) -- (replacement);
    \end{tikzpicture}
    \end{center}
    \vspace{0.08cm}
    \begin{questionbox}{Exemplo simples}
        Para o cromossomo \texttt{1 0 1 1 0}, o OneMax vale 3. Medimos o ciclo acima
        repetido 100 vezes, sem encerrar antes do fim.
    \end{questionbox}
\end{frame}

\begin{frame}{O que o experimento compara}
    \renewcommand{\arraystretch}{1.32}
    \begin{tabularx}{\textwidth}{>{\bfseries}p{0.23\textwidth} X}
        População $N$ & $128$, $512$ e $2048$ indivíduos. \\
        \hline
        Cromossomo $L$ & $64$, $256$ e $1024$ genes por indivíduo. \\
        \hline
        Tratamentos & Grade fatorial $3 \times 3$: todas as nove combinações de $N$ e $L$. \\
        \hline
        Controles & C++20, uma thread, OneMax, torneio 2, crossover de um ponto, $p_c=0{,}90$, mutação bit-flip, $p_m=0{,}01$ e $G=100$. \\
    \end{tabularx}
    \vspace{0.28cm}
    \begin{questionbox}{Por que uma grade fatorial?}
        Ela mostra o efeito de aumentar a população, o cromossomo e os dois juntos,
        sem trocar as regras do AG em cada caso.
    \end{questionbox}
\end{frame}

\begin{frame}{Como uma campanha produz comparações justas}
    \begin{columns}[T,onlytextwidth]
        \begin{column}{0.48\textwidth}
            \textbf{Antes e durante a coleta}
            \begin{itemize}
                \item afinidade no CPU lógico 0;
                \item aquecimento fora da medição;
                \item energia AC confirmada;
                \item mesmo binário e mesmos controles.
            \end{itemize}
        \end{column}
        \begin{column}{0.48\textwidth}
            \textbf{Planejamento das execuções}
            \begin{itemize}
                \item ordem aleatória no CSV;
                \item sementes registradas;
                \item 10 réplicas de tempo por tratamento;
                \item 5 réplicas para fases, IPC e cache.
            \end{itemize}
        \end{column}
    \end{columns}
    \vspace{0.25cm}
    \begin{questionbox}{Lote $B$}
        Algumas configurações são muito rápidas. Nelas, uma observação mede $B$ execuções
        completas do mesmo AG para formar uma duração estável. $B$ foi calibrado por tratamento,
        sem mudar $N$, $L$ ou $G$.
    \end{questionbox}
\end{frame}

\begin{frame}{Métricas: o que medimos e por quê}
    \small
    \renewcommand{\arraystretch}{1.16}
    \begin{tabularx}{\textwidth}{>{\bfseries}p{0.22\textwidth} p{0.27\textwidth} X}
        Medida & Significado simples & Pergunta que responde \\
        \hline
        $T_g=T/(BG)$ & tempo por geração & Quanto custa uma geração do AG? \\
        \hline
        $Q=BNLG/T$ & genes processados por segundo & O custo por gene muda entre os tratamentos? \\
        \hline
        Fases & parcela do tempo em cada fase do AG & Em que parte do AG o tempo se concentra? \\
        \hline
        IPC e cache & instruções por ciclo e taxa de faltas de cache & Há pistas de mudança no comportamento do processador? \\
    \end{tabularx}
    \vspace{0.14cm}
    \footnotesize\textcolor{alertred}{\textbf{Cuidado.}} Tempo, IPC e cache vieram de campanhas
    separadas. Os contadores ajudam a levantar hipóteses, mas não demonstram uma causa sozinhos.
\end{frame}

\begin{frame}{Três modos de medição, uma função para cada um}
    \begin{columns}[T,onlytextwidth]
        \begin{column}{0.31\textwidth}
            \begin{questionbox}{\texttt{baseline}}
                Mede o laço completo de gerações. É a fonte principal de $T_g$ e $Q$.
            \end{questionbox}
        \end{column}
        \begin{column}{0.31\textwidth}
            \begin{questionbox}{\texttt{phases}}
                Mede as cinco fases internas. Serve para explicar a composição do tempo.
            \end{questionbox}
        \end{column}
        \begin{column}{0.31\textwidth}
            \begin{questionbox}{\texttt{perf}}
                Coleta dois grupos: IPC e cache. Não mistura eventos de execuções diferentes.
            \end{questionbox}
        \end{column}
    \end{columns}
    \vspace{0.32cm}
    \begin{impactbox}{Regra de validade dos contadores}
        IPC usa ciclos e instruções da mesma linha \texttt{ipc\_user}. Cache usa referências
        e faltas da mesma linha \texttt{cache\_user}. Todos os grupos aceitos tiveram 100\% de tempo ativo.
    \end{impactbox}
\end{frame}

\begin{frame}{Ambiente observado e limite da comparação}
    \begin{columns}[T,onlytextwidth]
        \begin{column}{0.49\textwidth}
            \textbf{Plataforma usada}
            \begin{itemize}
                \item Intel Core i7-1355U, Linux 7.0.0-42;
                \item GCC 13.3.0 e arquitetura x86\_64;
                \item CPU lógico 0, com afinidade confirmada;
                \item L1d: 352 KiB, L2: 6,5 MiB, L3: 12 MiB.
            \end{itemize}
        \end{column}
        \begin{column}{0.48\textwidth}
            \textbf{Metadados preservados}
            \begin{itemize}
                \item versão do kernel e das ferramentas;
                \item topologia e caches da CPU;
                \item governor, Turbo, memória e alimentação;
                \item hash do binário e controles do AG.
            \end{itemize}
        \end{column}
    \end{columns}
    \vspace{0.22cm}
    \begin{impactbox}{Limite}
        Os resultados descrevem esta plataforma. As capacidades de cache ajudam a dar contexto,
        porém não identificam sozinhas uma causa microarquitetural.
    \end{impactbox}
\end{frame}

\begin{frame}{Da configuração ao gráfico, sem edição manual}
    \begin{center}
    \begin{tikzpicture}[node distance=0.35cm]
        \node[flowbox, text width=2.35cm] (config) {\textbf{Código e controles}\\AG e configuração};
        \node[flowbox, text width=2.35cm, right=of config] (design) {\textbf{DOE em CSV}\\ordem, sementes, lotes e réplicas};
        \node[flowbox, text width=2.35cm, right=of design] (raw) {\textbf{Dados brutos}\\comando, saída, status e ambiente};
        \node[flowbox, text width=2.35cm, right=of raw] (report) {\textbf{Validação e análise}\\CSV tratado, HTML e SVGs};
        \draw[flowarrow] (config) -- (design);
        \draw[flowarrow] (design) -- (raw);
        \draw[flowarrow] (raw) -- (report);
    \end{tikzpicture}
    \end{center}
    \vspace{0.28cm}
    \begin{questionbox}{Por que essa cadeia importa?}
        Cada resultado pode voltar ao \texttt{run\_id}, à saída original, ao binário e ao
        desenho que planejou a execução. O caderno Quarto gera as figuras a partir dos dados validados.
    \end{questionbox}
\end{frame}

\begin{frame}{Dados validados e como ler os gráficos}
    \begin{columns}[T,onlytextwidth]
        \begin{column}{0.48\textwidth}
            \textbf{Checagens antes da análise}
            \begin{itemize}
                \item esquema, status e tempo positivo;
                \item núcleo e controles congelados;
                \item contratos de fases e contadores.
            \end{itemize}
        \end{column}
        \begin{column}{0.48\textwidth}
            \textbf{Dados aceitos}
            \begin{itemize}
                \item 90 observações \texttt{baseline};
                \item 45 observações \texttt{phases};
                \item 45 para IPC e 45 para cache, sem remoções.
            \end{itemize}
        \end{column}
    \end{columns}
    \vspace{0.16cm}
    \begin{questionbox}{Chave para os gráficos de pontos}
        \textbf{Ponto} = mediana das réplicas de um tratamento. \textbf{Barra} = intervalo bootstrap
        de 95\% para essa mediana: reamostramos as réplicas 5.000 vezes, recalculamos a mediana e
        usamos os percentis de 2,5\% e 97,5\%. A barra não mostra o mínimo e o máximo das réplicas.
    \end{questionbox}
\end{frame}

\begin{frame}{Quanto custa uma geração do AG?}
    \chartimage{../figures/tempo-geracao-heatmap}
    \chartcaption{Cada célula combina $N$ e $L$. A mediana de $T_g$ cresce de 56,3~$\mu$s para
    13,2 ms entre os extremos da grade.}
\end{frame}

\begin{frame}{Vazão por gene: pontos com intervalo}
    \chartimage{../figures/vazao-workset}
    \chartcaption{\textbf{Eixos:} cada painel fixa $L$, o eixo $x$ mostra $N$ e o eixo $y$ mostra $Q$.
    \textbf{Ponto:} mediana de $Q=BNLG/T$ em 10 réplicas \texttt{baseline}. \textbf{Barra:} IC bootstrap
    de 95\%. \textbf{Uso:} comparar a velocidade por gene entre os tratamentos.}
\end{frame}

\begin{frame}{Em que parte do AG o tempo se concentra?}
    \chartimage{../figures/composicao-fases}
    \chartcaption{Cada barra soma 100\% do tempo médio das cinco réplicas \texttt{phases}. A mutação
    ocupa a maior parcela em todos os tratamentos.}
\end{frame}

\begin{frame}{Overhead da instrumentação: pontos horizontais}
    \chartimage{../figures/overhead-fases}
    \chartcaption{\textbf{Eixos:} cada linha é um tratamento, e o eixo horizontal mostra a diferença percentual.
    \textbf{Ponto:} mediana de $100(T_{phases}-T_{baseline})/T_{baseline}$ em 5 pares da mesma configuração.
    \textbf{Barra:} IC bootstrap de 95\%. A linha em 0\% indica nenhuma diferença. \textbf{Uso:} medir o efeito
    dos temporizadores no tempo principal.}
\end{frame}

\begin{frame}{Instruções por ciclo (IPC): pontos com intervalo}
    \chartimage{../figures/ipc-workset}
    \chartcaption{\textbf{Eixos:} cada painel fixa $L$, o eixo $x$ mostra $N$ e o eixo $y$ mostra IPC.
    \textbf{Ponto:} mediana de $instructions/cycles$ em 5 réplicas \texttt{ipc\_user}. \textbf{Barra:}
    IC bootstrap de 95\%. \textbf{Uso:} verificar se o ritmo de instruções por ciclo muda entre tratamentos.}
\end{frame}

\begin{frame}{Taxa de faltas de cache: pontos com intervalo}
    \chartimage{../figures/cache-miss-workset}
    \chartcaption{\textbf{Eixos:} cada painel fixa $L$, o eixo $x$ mostra $N$ e o eixo $y$ mostra a taxa.
    \textbf{Ponto:} mediana de $100\,cache\_misses/cache\_references$ em 5 réplicas \texttt{cache\_user}.
    \textbf{Barra:} IC bootstrap de 95\%. \textbf{Uso:} procurar evidência compatível com pressão de memória.}
\end{frame}

\begin{frame}{Variação entre réplicas: pontos horizontais}
    \chartimage{../figures/variabilidade-baseline}
    \chartcaption{\textbf{Eixos:} cada linha é um tratamento e o eixo horizontal mostra o coeficiente de
    variação. \textbf{Ponto:} $100\,desvio\text{-}padrão/média$ de $T_g$ nas 10 réplicas \texttt{baseline}.
    Não há barra porque este gráfico resume dispersão observada. \textbf{Uso:} avaliar se as medianas são estáveis.}
\end{frame}

\begin{frame}{Dificuldades que levaram a decisões explícitas}
    \renewcommand{\arraystretch}{1.27}
    \begin{tabularx}{\textwidth}{>{\bfseries}p{0.27\textwidth} X X}
        Evidência & Decisão tomada & Limite atual \\
        \hline
        Na grade 4 por 4, o maior caso levou cerca de 22 s. & Reduzimos a grade para 3 por 3 antes da campanha e mantivemos $G=100$. & O estudo não cobre o extremo antigo. \\
        \hline
        O primeiro conjunto de \texttt{perf} teve 91\% de tempo ativo. & Separamos IPC e cache e fixamos o \texttt{perf} no CPU 0. & Não combinamos grupos como uma única observação. \\
        \hline
        Os temporizadores internos mudam a execução. & Medimos \texttt{phases} e comparamos pares com \texttt{baseline}. & O overhead permanece um diagnóstico entre campanhas. \\
    \end{tabularx}
\end{frame}

\begin{frame}{Cobertura dos critérios da Etapa 2: método e infraestrutura}
    \renewcommand{\arraystretch}{1.28}
    \begin{tabularx}{\textwidth}{>{\bfseries}p{0.20\textwidth} X >{\bfseries}p{0.20\textwidth} X}
        DESC & ambiente real, CPU e metadados & DOE & grade fatorial, CSV, sementes e lotes \\
        \hline
        REPRO & código, hashes, dados brutos e \texttt{make all} & BOAS & afinidade, aquecimento, réplicas e ordem aleatória \\
        \hline
        DSBOX & dados tidy, Quarto, Python e geração automática & & \\
    \end{tabularx}
    \vspace{0.30cm}
    \begin{questionbox}{Evidência que pode ser conferida}
        \texttt{docs/environment.md}, \texttt{config/experiment.yaml}, \texttt{config/design.csv},
        dados em \texttt{data/raw/} e scripts no repositório registram esses critérios.
    \end{questionbox}
\end{frame}

\begin{frame}{Conclusão: andamento do experimento}
    \begin{columns}[T,onlytextwidth]
        \begin{column}{0.49\textwidth}
            \begin{questionbox}{O que já foi realizado}
                \begin{itemize}
                    \item AG implementado e validado com testes determinísticos;
                    \item planejamento fatorial $3 \times 3$ executado;
                    \item campanhas de tempo, fases, IPC e cache coletadas;
                    \item dados validados e rastreáveis até cada execução;
                    \item caderno de análise e sete visualizações gerados.
                \end{itemize}
            \end{questionbox}
        \end{column}
        \begin{column}{0.47\textwidth}
            \begin{questionbox}{O que os dados indicam}
                \begin{itemize}
                    \item o tempo por geração aumenta com $N$ e $L$;
                    \item a vazão normalizada por gene permanece em faixa próxima;
                    \item mutação concentra a maior parcela do tempo;
                    \item IPC e faltas de cache oferecem evidências descritivas,
                    não uma explicação causal.
                \end{itemize}
            \end{questionbox}
        \end{column}
    \end{columns}
    \vspace{0.22cm}
    \begin{impactbox}{Estado atual e continuidade}
        A Etapa 2 está consolidada em uma campanha reproduzível, com resultados,
        limitações e decisões documentados. O próximo passo é aprofundar a análise
        estatística, investigar as hipóteses levantadas e ampliar os experimentos
        somente quando houver justificativa.
    \end{impactbox}
\end{frame}

\end{document}
